\documentclass[aspectratio=169, 10pt]{beamer}

% --- Packages ---
\usepackage[utf8]{inputenc}
\usepackage[vietnamese]{babel}
\usepackage{amsmath, amssymb, amsfonts}
\usepackage{booktabs}
\usepackage{tikz}
\usetikzlibrary{shapes.geometric, arrows.meta, positioning, calc, fit, backgrounds, shadows}
\usepackage{xcolor}
\usepackage{colortbl}
\usepackage{listings}
\usepackage{pifont}

% --- Beamer Theme & Styling ---
\usetheme{Madrid}
\usecolortheme{whale}

% Định nghĩa màu sắc thương hiệu học thuật hiện đại
\definecolor{primaryBlue}{RGB}{24, 76, 120}
\definecolor{accentOrange}{RGB}{230, 81, 0}
\definecolor{darkSlate}{RGB}{38, 50, 56}
\definecolor{lightGray}{RGB}{245, 247, 250}
\definecolor{successGreen}{RGB}{46, 125, 50}
\definecolor{dangerRed}{RGB}{198, 40, 40}
\definecolor{purpleTag}{RGB}{123, 31, 162}

\setbeamercolor{palette primary}{bg=primaryBlue, fg=white}
\setbeamercolor{palette secondary}{bg=primaryBlue!85, fg=white}
\setbeamercolor{palette tertiary}{bg=primaryBlue!70, fg=white}
\setbeamercolor{structure}{fg=primaryBlue}
\setbeamercolor{title}{fg=white, bg=primaryBlue}
\setbeamercolor{frametitle}{fg=white, bg=primaryBlue}
\setbeamercolor{block title}{bg=primaryBlue!15, fg=primaryBlue}
\setbeamercolor{block body}{bg=lightGray, fg=darkSlate}
\setbeamercolor{block title alerted}{bg=dangerRed!15, fg=dangerRed}
\setbeamercolor{block body alerted}{bg=dangerRed!5, fg=darkSlate}
\setbeamercolor{block title example}{bg=successGreen!15, fg=successGreen}
\setbeamercolor{block body example}{bg=successGreen!5, fg=darkSlate}

\setbeamertemplate{navigation symbols}{}
\setbeamertemplate{footline}[frame number]

% --- Metadata ---
\title[\textbf{Báo Cáo Tiến Độ Tuần 2: Dual-Agent Reflexion}]{\textbf{Bước Chuyển Mô Thức: Kiến Trúc Dual-Agent Reflexion \& Công Cụ Thu Phóng Thời Gian Liên Tục Trong Video Agent}}
\subtitle{Báo Cáo Tiến Độ Nghiên Cứu Tuần 2: Tái Cấu Trúc Toàn Diện Hệ Thống Agentic Video Understanding}
\author[Lê Nhật Sơn]{Học viên: \textbf{Lê Nhật Sơn}}
\institute[Computer Vision Lab]{
    Nghiên cứu về Long-Horizon Video Understanding với Multimodal Agents \\
    Đột phá từ nền tảng VideoSeek (CVPR 2026) sang Kiến trúc Phản Tư Đa Tác Tử
}
\date{\today}

\begin{document}

% -----------------------------------------------------------------------------
% SLIDE 1: TITLE PAGE
% -----------------------------------------------------------------------------
\begin{frame}
    \titlepage
\end{frame}

% -----------------------------------------------------------------------------
% SLIDE 2: ĐẶT VẤN ĐỀ & NHẬN THỨC MỚI
% -----------------------------------------------------------------------------
\begin{frame}{1. Đặt Vấn Đề \& Nhận Thức Mới: Tại Sao 3 Công Cụ Rời Rạc Là Chưa Đủ?}
    \small
    \begin{columns}[T]
        \begin{column}{0.32\textwidth}
            \begin{alertblock}{\textbf{1. Bản Chất Toán Học}}
                \begin{itemize}\footnotesize
                    \item 3 công cụ gốc (\texttt{overview}, \texttt{skim}, \texttt{focus}) thực chất chỉ là 3 điểm rời rạc trên cùng hàm:
                    $$\text{Sample}(start, end, \text{fps})$$
                    \item Việc chia nhỏ tạo ra ranh giới nhân tạo cứng nhắc.
                \end{itemize}
            \end{alertblock}
        \end{column}

        \begin{column}{0.32\textwidth}
            \begin{alertblock}{\textbf{2. Ranh Giới Gây Bối Rối}}
                \begin{itemize}\footnotesize
                    \item Cần xem đoạn $14$s: dùng \texttt{focus} thì bị khóa cứng $1$ FPS (chỉ 14 frame); dùng \texttt{skim} thì vi phạm luật $>16$s.
                    \item Agent tốn quá nhiều turn suy nghĩ "chọn tool nào" thay vì tập trung vào mục tiêu.
                \end{itemize}
            \end{alertblock}
        \end{column}

        \begin{column}{0.32\textwidth}
            \begin{alertblock}{\textbf{3. Bẫy Lặp Phân Vân}}
                \begin{itemize}\footnotesize
                    \item Khi đã thấy bằng chứng (ví dụ: két sắt tại Turn 8), cơ chế ReAct thuần túy không có bộ phận tự vấn.
                    \item Dẫn tới bẫy lặp phân vân 12 turns (kéo dài đến Turn 20 fallback).
                \end{itemize}
            \end{alertblock}
        \end{column}
    \end{columns}

    \vspace{0.25cm}
    \begin{exampleblock}{\textbf{Bước Chuyển Mô Thức Đề Xuất (Paradigm Shift)}}
        \centering
        \small Thay thế 3 tools rời rạc bằng \textbf{01 Công cụ Thu phóng Liên tục duy nhất} (\texttt{inspect}) kết hợp \textbf{Kiến trúc Dual-Agent Reflexion} (Actor Agent + Reflexion Evaluator Agent).
    \end{exampleblock}
\end{frame}

% -----------------------------------------------------------------------------
% SLIDE 3: SƠ ĐỒ KIẾN TRÚC HỆ THỐNG MỚI (TIKZ DIAGRAM)
% -----------------------------------------------------------------------------
\begin{frame}{2. Sơ Đồ Kiến Trúc Hệ Thống: Dual-Agent Reflexion \& Unified Tool}
\centering
\resizebox{0.92\textwidth}{!}{%
\begin{tikzpicture}[
    font=\small,
    node distance=1.1cm and 1.3cm,
    base/.style={draw, rounded corners=4pt, align=center, thick, drop shadow},
    inputNode/.style={base, fill=primaryBlue!10, draw=primaryBlue, minimum width=2.6cm, minimum height=0.9cm},
    actorNode/.style={base, fill=primaryBlue!20, draw=primaryBlue, minimum width=3.6cm, minimum height=1.1cm, font=\bfseries\small},
    toolNode/.style={base, fill=accentOrange!20, draw=accentOrange, minimum width=3.8cm, minimum height=1.1cm, font=\bfseries\small},
    obsNode/.style={base, fill=lightGray, draw=darkSlate, minimum width=3.2cm, minimum height=0.9cm},
    evalNode/.style={base, fill=purpleTag!20, draw=purpleTag, minimum width=4.0cm, minimum height=1.1cm, font=\bfseries\small},
    decisionNode/.style={base, diamond, aspect=2, fill=accentOrange!15, draw=accentOrange, inner sep=2pt},
    answerNode/.style={base, fill=successGreen, draw=successGreen!80!black, text=white, minimum width=2.6cm, minimum height=0.9cm, font=\bfseries},
    arrow/.style={-{Stealth[length=3mm]}, thick, draw=darkSlate},
    dashArrow/.style={-{Stealth[length=3mm]}, dashed, thick, draw=purpleTag}
]

    % --- Nodes ---
    \node[inputNode] (input) {\textbf{Input Video $V$} \\ + \textbf{User Query $Q$}};
    \node[actorNode, right=1.0cm of input] (actor) {\textbf{1. ACTOR AGENT} \\ \footnotesize System Prompt: \texttt{SYSTEM\_PROMPT} \\ \scriptsize Sinh Thought \& Chọn $(start, end, \text{fps})$};
    
    \node[toolNode, right=1.2cm of actor] (tool) {\textbf{2. UNIFIED TOOL: \texttt{inspect}} \\ \scriptsize $N = (end - start) \times \text{fps} \le 64$ frames \\ \scriptsize Adaptive: Macro Grid ($N>16$) vs Micro ($N\le 16$)};
    
    \node[obsNode, right=1.0cm of tool] (obs) {\textbf{Observation $o_t$} \\ \footnotesize Khung hình + Subtitles + Timestamps};
    
    \node[evalNode, below=1.2cm of obs] (eval) {\textbf{3. REFLEXION EVALUATOR} \\ \footnotesize System Prompt: \texttt{EVALUATOR\_PROMPT} \\ \scriptsize Thẩm định độc lập: Sufficiency \& Loop Check};
    
    \node[decisionNode, below=1.2cm of actor] (check) {Bằng chứng đủ?};
    \node[answerNode, left=1.2cm of check] (answer) {\textbf{\texttt{<answer>}} \\ Chốt Đáp Án};

    % --- Edges ---
    \draw[arrow] (input) -- (actor);
    \draw[arrow] (actor) -- node[above, font=\scriptsize] {Gọi tool} (tool);
    \draw[arrow] (tool) -- (obs);
    \draw[arrow] (obs) -- node[right, font=\scriptsize] {Gửi quan sát} (eval);
    \draw[arrow] (eval) -- (check);
    
    \draw[arrow] (check) -- node[above, font=\scriptsize, text=successGreen] {\textbf{YES (SUFFICIENT)}} (answer);
    
    % Feedback loop phản tư
    \draw[dashArrow] (check) -- node[left, font=\scriptsize, text=dangerRed] {\textbf{NO (INSUFFICIENT)}} (actor)
        node[pos=0.5, right=0.1cm, font=\scriptsize, text=purpleTag] {Chỉ thị phản tư: Điều chỉnh $[start, end]$ \& Tăng FPS};

    % Khung bao quanh cơ chế mới
    \begin{scope}[on background layer]
        \node[draw=purpleTag, dashed, fill=purpleTag!5, rounded corners=6pt, fit=(actor) (tool) (eval) (check), inner sep=0.35cm] (boxReflexion) {};
        \node[above=0.05cm of boxReflexion.north, text=purpleTag, font=\bfseries\footnotesize] {Kiến Trúc Dual-Agent Reflexion \& Continuous Temporal Zooming};
    \end{scope}

\end{tikzpicture}
}
\vspace{0.2cm}
\begin{itemize}\footnotesize
    \item \textbf{Actor Agent}: Tập trung ra quyết định và thu phóng thời gian qua duy nhất 1 tool \texttt{inspect}.
    \item \textbf{Reflexion Evaluator Agent}: Đóng vai trò Giám khảo phản biện độc lập, ngắt bẫy lặp ngay khi đủ bằng chứng.
\end{itemize}
\end{frame}

% -----------------------------------------------------------------------------
% SLIDE 4: CHI TIẾT CÔNG CỤ THU PHÓNG LIÊN TỤC (INSPECT)
% -----------------------------------------------------------------------------
\begin{frame}{3. Chi Tiết Kỹ Thuật: Unified Continuous Inspection Tool (\texttt{inspect})}
    \begin{columns}[T]
        \begin{column}{0.48\textwidth}
            \begin{block}{\textbf{1. Khống Chế Ngân Sách Frame (2 Lớp)}}
                \begin{itemize}
                    \item \textbf{Lớp 1 (Cognitive Guardrail)}:
                    \begin{itemize}\footnotesize
                        \item Nêu rõ công thức trong System Prompt của cả 2 Agent:
                        $$(t_{\text{end}} - t_{\text{start}}) \times \text{fps} \le 64$$
                    \end{itemize}
                    \item \textbf{Lớp 2 (Hard Deterministic Guardrail)}:
                    \begin{itemize}\footnotesize
                        \item Trong code \texttt{inspect.py}, nếu LLM tính nhẩm sai vượt $64$ frames:
                        $$N_{\text{frames}} = \min(N_{\text{desired}}, 64)$$
                        \item Tự động gắn thông báo điều chỉnh $\implies$ \textbf{100\% không bao giờ tràn context hay sập API}.
                    \end{itemize}
                \end{itemize}
            \end{block}
        \end{column}

        \begin{column}{0.48\textwidth}
            \begin{block}{\textbf{2. Đóng Gói Ảnh Thích Ứng (Adaptive Packing)}}
                \begin{itemize}
                    \item \textbf{Macro Grid Mode ($N > 16$ frames)}:
                    \begin{itemize}\footnotesize
                        \item Dùng cho quét toàn cảnh / đoạn dài.
                        \item Làm tròn bội số của 8, ghép mỗi 8 frames thành lưới $2\times 4$ ($256$px) kèm \texttt{Timestamp Matrix}.
                        \item Tiết kiệm token tối đa khi quét vĩ mô.
                    \end{itemize}
                    \item \textbf{Micro High-Res Mode ($N \le 16$ frames)}:
                    \begin{itemize}\footnotesize
                        \item Dùng khi soi cận cảnh, kiểm chứng chi tiết.
                        \item Giữ nguyên từng frame riêng lẻ độ nét cao.
                        \item Timestamp chính xác đến phần trăm giây \texttt{\{t:.2f\}s} khi $\text{fps} \ge 4.0$.
                    \end{itemize}
                \end{itemize}
            \end{block}
        \end{column}
    \end{columns}
\end{frame}

% -----------------------------------------------------------------------------
% SLIDE 5: CƠ CHẾ DUAL-AGENT & LOOP BREAKER
% -----------------------------------------------------------------------------
\begin{frame}{4. Kiến Trúc Dual-Agent: Tách Biệt Actor \& Reflexion Evaluator}
    \begin{columns}[T]
        \begin{column}{0.48\textwidth}
            \begin{block}{\textbf{1. Actor Agent (\texttt{VideoSeekAgent})}}
                \begin{itemize}
                    \item Sở hữu \texttt{SYSTEM\_PROMPT} riêng bám sát phong cách chuẩn của VideoSeek.
                    \item \textbf{Khởi đầu linh hoạt (Cold-Start Strategy)}:
                    \begin{itemize}\footnotesize
                        \item Nếu câu hỏi chưa có mốc thời gian $\implies$ Quét vĩ mô toàn cảnh với FPS thấp: $\text{fps} \approx 64 / T$.
                        \item Nếu câu hỏi/phụ đề đã rõ mốc thời gian $\implies$ Nhảy thẳng vào soi cận cảnh, \textbf{không bị ép gọi overview mù quáng ở Step 0}.
                    \end{itemize}
                    \item Tiếp thu phản hồi phản tư từ Evaluator để điều chỉnh $[start, end]$ và $\text{fps}$ ở bước kế tiếp.
                \end{itemize}
            \end{block}
        \end{column}

        \begin{column}{0.48\textwidth}
            \begin{exampleblock}{\textbf{2. Evaluator Agent (\texttt{ReflexionEvaluator})}}
                \begin{itemize}
                    \item Sở hữu \texttt{EVALUATOR\_SYSTEM\_PROMPT} và message context hoàn toàn độc lập.
                    \item \textbf{Sufficiency Verification}: Xác thực trực tiếp xem dữ kiện đã trả lời trọn vẹn câu hỏi hay chưa.
                    \item \textbf{Dập tắt bẫy lặp (Loop Breaker)}:
                    \begin{itemize}\footnotesize
                        \item Khi thấy két sắt sau ảnh Hitler (Turn 8 ở baseline cũ) $\implies$ Đánh giá ngay \textbf{\texttt{SUFFICIENT}} $\implies$ Ra lệnh gọi \texttt{answer} dừng ngay lập tức.
                        \item Nếu \textbf{\texttt{INSUFFICIENT}} $\implies$ Chỉ rõ thiếu sót và định hướng tăng FPS hoặc dịch chuyển cửa sổ thời gian.
                    \end{itemize}
                \end{itemize}
            \end{exampleblock}
        \end{column}
    \end{columns}
\end{frame}

% -----------------------------------------------------------------------------
% SLIDE 6: KẾT QUẢ KIỂM THỬ TỰ ĐỘNG
% -----------------------------------------------------------------------------
\begin{frame}{5. Kết Quả Kiểm Thử Độc Lập: 8/8 Unit Tests Đạt Trạng Thái PASSED}
    \begin{table}[h]
        \centering
        \resizebox{0.95\textwidth}{!}{%
        \begin{tabular}{lllc}
            \toprule
            \textbf{STT} & \textbf{Kịch Bản Kiểm Thử (\texttt{test\_unified\_reflexion.py})} & \textbf{Mục Tiêu Đánh Giá} & \textbf{Kết Quả} \\
            \midrule
            1 & \texttt{test\_1\_inspect\_tool\_schema} & Schema nghiêm ngặt (\texttt{strict=True}), đủ 4 tham số bắt buộc & \textcolor{successGreen}{\textbf{PASSED}} \\
            2 & \texttt{test\_2\_execute\_inspect\_macro\_grid} & Đóng gói lưới $2\times 4$ khi $N > 16$ kèm Timestamp Matrix & \textcolor{successGreen}{\textbf{PASSED}} \\
            3 & \texttt{test\_3\_execute\_inspect\_micro\_highres} & Ảnh đơn độ nét cao khi $N \le 16$, timestamp phần trăm giây \texttt{.2f} & \textcolor{successGreen}{\textbf{PASSED}} \\
            4 & \texttt{test\_4\_frame\_budget\_clamping} & Tự động clamp an toàn khi yêu cầu vượt trần ($480 \implies 64$ frames) & \textcolor{successGreen}{\textbf{PASSED}} \\
            5 & \texttt{test\_5\_dual\_agent\_prompt\_separation} & Tách biệt System Prompt \& context giữa Actor và Evaluator & \textcolor{successGreen}{\textbf{PASSED}} \\
            6 & \texttt{test\_6\_end\_to\_end\_reflexion\_workflow} & Luồng phối hợp hoàn chỉnh: Zoom $\implies$ Sufficiency $\implies$ Chốt đáp án & \textcolor{successGreen}{\textbf{PASSED}} \\
            7 & \texttt{test\_7\_inspect\_boundary\_parameters} & Xử lý biên: $t_{\text{start}} \ge t_{\text{end}}$, thời gian âm, vượt duration, $\text{fps} \le 0$ & \textcolor{successGreen}{\textbf{PASSED}} \\
            8 & \texttt{test\_8\_pydantic\_v2\_and\_baseline\_3tools} & Tương thích Pydantic v2 tool calls \& luồng Baseline đối chứng & \textcolor{successGreen}{\textbf{PASSED}} \\
            \bottomrule
        \end{tabular}%
        }
    \end{table}

    \vspace{-0.1cm}
    \begin{block}{\textbf{Đánh giá chất lượng kỹ thuật:}}
        \centering
        \footnotesize Toàn bộ 8 bài test chạy tự động trong \textbf{0.315 giây}. Mã nguồn đạt độ ổn định cao, không có nguy cơ rò rỉ bộ nhớ hay lỗi cú pháp.
    \end{block}
\end{frame}

% -----------------------------------------------------------------------------
% SLIDE 7: A/B BENCHMARK CELL 7 TRÊN COLAB
% -----------------------------------------------------------------------------
\begin{frame}{6. Cơ Chế Kiểm Thử Đối Đầu Trực Diện A/B Benchmark (Cell 7)}
    \begin{table}[h]
        \centering
        \resizebox{0.95\textwidth}{!}{%
        \begin{tabular}{lcc}
            \toprule
            \textbf{Tiêu Chí Đánh Giá} & \textbf{BASELINE (VideoSeek Gốc)} & \textbf{UPGRADED (Kiến Trúc Mới Đề Xuất)} \\
            \midrule
            \textbf{Kiến trúc Agent} & Single-Agent (Không có phản tư) & \textbf{Dual-Agent (Actor + Reflexion Evaluator)} \\
            \textbf{Bộ công cụ thị giác} & 3 tools cứng: \texttt{overview}, \texttt{skim}, \texttt{focus} & \textbf{01 Unified Tool duy nhất: \texttt{inspect(t1, t2, fps)}} \\
            \textbf{Cơ chế đóng gói ảnh} & Cố định cứng theo từng tool & \textbf{Linh hoạt: Macro 2x4 Grid / Micro High-Res} \\
            \textbf{Khả năng điều chỉnh FPS} & Khóa chết $1$ FPS tại \texttt{focus} & \textbf{Động liên tục ($0.02 - 8.0$ FPS)} \\
            \textbf{Khởi đầu lượt 1 (Cold-start)} & Luôn bị ép chạy \texttt{overview} & \textbf{Chủ động soi ngay nếu câu hỏi có mốc thời gian} \\
            \textbf{Cơ chế dập tắt bẫy lặp} & Không có (kẹt phân vân tới 20 turns) & \textbf{Reflexion Evaluator ra lệnh chốt ngay khi đủ} \\
            \bottomrule
        \end{tabular}%
        }
    \end{table}

    \vspace{-0.1cm}
    \begin{exampleblock}{\textbf{Tích hợp sẵn sàng trên Notebook Colab (\texttt{test.ipynb}):}}
        \begin{itemize}\footnotesize
            \item \textbf{Cell 7}: Tự động thực thi đối đầu 2 lượt trên cùng 1 video \& 1 câu hỏi, tự động dò tìm file phụ đề \texttt{.srt}.
            \item In bảng đo lường trực quan (số turns, thời gian runtime, chuỗi tool gọi) và lưu vào file \texttt{outputs/ab\_comparison\_report.json}.
        \end{itemize}
    \end{exampleblock}
\end{frame}

% -----------------------------------------------------------------------------
% SLIDE 8: BƯỚC CHUYỂN MÔ THỨC & Ý NGHĨA KHOA HỌC
% -----------------------------------------------------------------------------
\begin{frame}{7. Bước Chuyển Mô Thức \& Giá Trị Đóng Góp Khoa Học}
    \begin{columns}[T]
        \begin{column}{0.48\textwidth}
            \begin{block}{\textbf{Tiến Trình Phát Triển Của Đề Tài}}
                \begin{itemize}\footnotesize
                    \item \textbf{Giai đoạn 1 (VideoSeek gốc)}:
                    3 tools rời rạc, $1$ FPS cứng, mù sự kiện đột ngột, dễ kẹt lặp.
                    \item \textbf{Giai đoạn 2 (Module 2 trước đó)}:
                    Thêm Query Router \& Dynamic FPS $\implies$ Đóng góp mang tính \textit{incremental} (chắp vá tham số).
                    \item \textbf{Giai đoạn 3 (Tuần này - Đột phá)}:
                    Tái cấu trúc thành \textbf{Continuous Temporal Inspection + Dual-Agent Reflexion} $\implies$ Bước chuyển mô thức thực thụ!
                \end{itemize}
            \end{block}
        \end{column}

        \begin{column}{0.48\textwidth}
            \begin{exampleblock}{\textbf{2 Đóng Góp Lớn Cho Bài Báo Khoa Học}}
                \begin{itemize}\footnotesize
                    \item \textbf{Đóng góp 1: Continuous Temporal Inspection Tool (CTIT)}:
                    Xóa bỏ ranh giới nhân tạo giữa các công cụ thị giác; biến bài toán quan sát video thành cơ chế điều khiển thu phóng thời gian liên tục $(t_{\text{start}}, t_{\text{end}}, \text{fps})$.
                    \item \textbf{Đóng góp 2: Reflexion-Augmented Video Decision Loop}:
                    Giải quyết tận gốc hạn chế kinh điển của ReAct (vòng lặp suy luận dư thừa và tê liệt quyết định) trong hiểu video dài.
                \end{itemize}
            \end{exampleblock}
        \end{column}
    \end{columns}
\end{frame}

% -----------------------------------------------------------------------------
% SLIDE 9: KẾ HOẠCH BƯỚC TIẾP THEO
% -----------------------------------------------------------------------------
\begin{frame}{8. Kế Hoạch Nghiên Cứu Tuần Tiếp Theo}
    \begin{columns}[T]
        \begin{column}{0.48\textwidth}
            \begin{alertblock}{\textbf{Trọng tâm 1: Benchmark Quy Mô Lớn}}
                \begin{itemize}
                    \item Chạy thực nghiệm đối đầu trên tập con phân tầng ($100-150$ videos) của \textbf{LVBench} và \textbf{VideoMME}.
                    \item Đo lường toàn diện 4 chỉ số:
                    \begin{itemize}\footnotesize
                        \item Độ chính xác (\%).
                        \item Số khung hình tiêu thụ (\#Frames).
                        \item Số lượt suy luận trung bình (\#Turns).
                        \item Chi phí \& Độ trễ API.
                    \end{itemize}
                \end{itemize}
            \end{alertblock}
        \end{column}

        \begin{column}{0.48\textwidth}
            \begin{exampleblock}{\textbf{Trọng tâm 2: Hoàn Thiện Bản Thảo Bài Báo}}
                \begin{itemize}
                    \item Hoàn thành bản thảo sơ bộ (Manuscript Draft):
                    \begin{itemize}\footnotesize
                        \item Mục \textit{Methodology}: Mô tả hình thức hóa toán học của Continuous Zooming và Dual-Agent Reflexion.
                        \item Mục \textit{Ablation Study}: Bóc tách định lượng đóng góp độc lập của Tool \texttt{inspect} vs Khối Reflexion Evaluator.
                    \end{itemize}
                \end{itemize}
            \end{exampleblock}
        \end{column}
    \end{columns}
\end{frame}

% -----------------------------------------------------------------------------
% SLIDE 10: KẾT LUẬN & Q&A
% -----------------------------------------------------------------------------
\begin{frame}
    \vspace{1.5cm}
    \centering
    \Huge \textbf{\textcolor{primaryBlue}{Cảm Ơn Thầy/Cô \& Các Bạn!}} \\
    \vspace{0.8cm}
    \normalsize \textit{Rất mong nhận được ý kiến đóng góp \& định hướng tiếp theo từ Thầy/Cô.} \\
    \vspace{0.5cm}
    \small
    Mã nguồn dự án: \url{https://github.com/nhatsonle/videoseek.git} \\
    \footnotesize (Nhánh \texttt{main}: Dual-Agent Reflexion | Nhánh \texttt{backup/module2-adaptive}: Phiên bản cũ)
\end{frame}

\end{document}
